\chapter{DF-001 --- Show DOGFOOD Coverage Progress}\label{record-064}

\noindent{\footnotesize\color{BookMuted}\path{docs/design/DF-001-show-dogfood-coverage-progress.md}\par}\smallskip

Legend: \href{https://github.com/flyingrobots/bijou/blob/1853ad902f6bf219326409e7f743a5a1d6517dbc/docs/legends/DF-dogfood-field-guide.md}{DF --- DOGFOOD Field Guide}

\section{Why this cycle exists}\label{record-064-why-this-cycle-exists}

We want to release Bijou without pretending DOGFOOD is already a complete reference.

The honest answer is not to block release forever. It is to make DOGFOOD visibly truthful about its current documentation coverage.

This cycle exists to put that truth directly into the initial docs experience:

\begin{itemize}
\tightlist
\item
  a progress indicator
\item
  an explicit ratio
\item
  a defensible denominator tied to the canonical component-family reference
\end{itemize}

\section{Scope of this cycle}\label{record-064-scope-of-this-cycle}

This cycle intentionally covers:

\begin{itemize}
\tightlist
\item
  a canonical DOGFOOD coverage metric based on documented component families
\item
  mapping current DOGFOOD stories onto that coverage model
\item
  a docs-pane progress bar and ratio
\item
  tests that prove the landing surface reflects real coverage instead of a hand-wavy claim
\end{itemize}

It does \textbf{not} include:

\begin{itemize}
\tightlist
\item
  expanding story coverage itself
\item
  claiming feature completeness
\item
  a broader release-management dashboard
\item
  automatic release gating
\end{itemize}

\section{Human users}\label{record-064-human-users}

\subsection{Primary human user}\label{record-064-primary-human-user}

A user opening DOGFOOD to decide how complete and trustworthy the field guide currently is.

They need:

\begin{itemize}
\tightlist
\item
  a clear signal that DOGFOOD is still in progress
\item
  an honest coverage ratio
\item
  enough context to understand that the metric is about documented component families
\end{itemize}

\subsection{Human hill}\label{record-064-human-hill}

A user can open DOGFOOD and immediately see that it is a living field guide with explicit progress, not a silently incomplete reference pretending to be finished.

\section{Agent users}\label{record-064-agent-users}

\subsection{Primary agent user}\label{record-064-primary-agent-user}

An agent maintaining DOGFOOD story coverage over time.

It needs:

\begin{itemize}
\tightlist
\item
  a stable coverage model
\item
  a defensible denominator
\item
  a way to add stories later and have the landing progress move automatically
\end{itemize}

\subsection{Agent hill}\label{record-064-agent-hill}

An agent can add or remove DOGFOOD story coverage and see the docs-pane progress update from real source data rather than manual copy edits.

\section{Human playback}\label{record-064-human-playback}

\begin{enumerate}
\def\labelenumi{\arabic{enumi}.}
\tightlist
\item
  A user launches DOGFOOD.
\item
  The landing screen shows the normal hero treatment.
\item
  The initial docs content pane shows a progress bar and ratio for current DOGFOOD documentation coverage.
\item
  The user understands that DOGFOOD is in progress and can see how much of the canonical family map is currently covered.
\end{enumerate}

\section{Agent playback}\label{record-064-agent-playback}

\begin{enumerate}
\def\labelenumi{\arabic{enumi}.}
\tightlist
\item
  An agent opens the DF-001 cycle doc.
\item
  It sees that the denominator is the canonical component-family reference, not a vague export count.
\item
  It runs the cycle and app tests.
\item
  It can verify that the docs-pane copy and progress bar are driven from real story coverage data.
\item
  It can point to the next backlog item for actually expanding coverage.
\end{enumerate}

\section{Linked invariants}\label{record-064-linked-invariants}

\begin{itemize}
\tightlist
\item
  \href{https://github.com/flyingrobots/bijou/blob/cbd70612140e7a450d764b3c7be1c6af878b99bc/docs/invariants/tests-are-the-spec.md}{Tests Are the Spec}
\item
  \href{https://github.com/flyingrobots/bijou/blob/cbd70612140e7a450d764b3c7be1c6af878b99bc/docs/invariants/visible-controls-are-a-promise.md}{Visible Controls Are a Promise}
\item
  \href{https://github.com/flyingrobots/bijou/blob/cbd70612140e7a450d764b3c7be1c6af878b99bc/docs/invariants/docs-are-the-demo.md}{Docs Are the Demo}
\end{itemize}

\section{Invariants for this cycle}\label{record-064-invariants-for-this-cycle}

\begin{itemize}
\tightlist
\item
  DOGFOOD must not imply fuller coverage than it actually has
\item
  the coverage denominator must be explicit and stable
\item
  adding story coverage later should update the progress metric automatically
\item
  the title screen must not present documentation coverage as though it were a loading state
\item
  the docs-pane coverage UI must remain secondary to the primary onboarding content
\end{itemize}

\section{Implementation outline}\label{record-064-implementation-outline}

\begin{enumerate}
\def\labelenumi{\arabic{enumi}.}
\tightlist
\item
  Create a DOGFOOD legend and coverage-focused cycle doc.
\item
  Define the canonical coverage denominator from the design-system component-family reference.
\item
  Map current DOGFOOD stories onto that family coverage model.
\item
  Write failing cycle and app tests for the new landing progress treatment.
\item
  Render a progress bar and ratio in the initial docs content pane.
\item
  Close the cycle and spawn the next DOGFOOD coverage backlog item.
\end{enumerate}

\section{Tests to write first}\label{record-064-tests-to-write-first}

Under \texttt{t\allowbreak{}e\allowbreak{}s\allowbreak{}t\allowbreak{}s\allowbreak{}/\allowbreak{}c\allowbreak{}y\allowbreak{}c\allowbreak{}l\allowbreak{}e\allowbreak{}s\allowbreak{}/\allowbreak{}D\allowbreak{}F\allowbreak{}-\allowbreak{}0\allowbreak{}0\allowbreak{}1\allowbreak{}/}:

\begin{itemize}
\tightlist
\item
  the active cycle doc includes the required workflow sections
\item
  DOGFOOD coverage resolves from canonical family reference data
\item
  the initial docs content surface shows the coverage ratio and percent
\item
  the next DOGFOOD backlog item exists
\end{itemize}

\section{Risks}\label{record-064-risks}

\begin{itemize}
\tightlist
\item
  choosing a denominator that feels arbitrary or magical
\item
  making the title screen feel like a loading screen when it is actually a title screen
\item
  hardcoding a number that will drift the moment a new story lands
\end{itemize}

\section{Out of scope}\label{record-064-out-of-scope}

\begin{itemize}
\tightlist
\item
  story authoring for additional families
\item
  docs completeness gating in CI
\item
  release automation
\end{itemize}

\section{Retrospective}\label{record-064-retrospective}

\subsection{What landed}\label{record-064-what-landed}

\begin{itemize}
\tightlist
\item
  a new DOGFOOD legend
\item
  a real coverage model based on canonical component families
\item
  story-to-family coverage mapping
\item
  a docs-pane progress bar and ratio that update from the current story set
\end{itemize}

\subsection{Drift from ideal}\label{record-064-drift-from-ideal}

This cycle does not improve DOGFOOD coverage itself.

That is intentional. The goal here is honesty first, not pretending that a progress indicator is the same thing as documentation work.

\subsection{Debt spawned}\label{record-064-debt-spawned}

Spawned:

\begin{itemize}
\tightlist
\item
  \href{https://github.com/flyingrobots/bijou/blob/cbd70612140e7a450d764b3c7be1c6af878b99bc/docs/method/graveyard/legacy-backlog/DF-002-expand-dogfood-component-family-coverage.md}{DF-002 --- Expand DOGFOOD Component Family Coverage}
\end{itemize}

